\section{Related work}
\label{sec:related}

\subsection{Indirect prompt injection and agent benchmarks}
Prompt injection was first shown against models that read attacker text in their prompt~\cite{perez2022ignore}; Greshake et al.~\cite{greshake2023} showed that the attacker does not need access to the prompt, because an application that retrieves web pages, emails, or documents places third-party text in the model's context. Liu et al.~\cite{liu2024formalizing} formalized the attack as a target task and an injected task and benchmarked attacks and defenses on that basis. Benchmarks for agents followed. InjecAgent~\cite{injecagent} places 62 attacker instructions inside simulated tool responses; AgentDojo~\cite{agentdojo} provides four suites of realistic tasks with stateful environments, ground-truth tool calls, and injection points in the data; BIPIA~\cite{bipia} injects attacks into emails, tables, code answers, and other external content; ASB~\cite{asb} covers ten agent scenarios and many attack and defense combinations; WASP~\cite{wasp} evaluates web agents in realistic, sandboxed web environments; and AgentDyn~\cite{agentdyn} adds open-ended tasks and benign third-party instructions that a defense must not suppress. $\tau$-bench~\cite{taubench} is not a security benchmark, but its ground-truth action sequences give a second, independently built agent environment. PIArena~\cite{piarena} integrates attacks, defenses, and benchmarks in one platform and reports that defenses generalize poorly across tasks. These benchmarks measure whether an agent or model follows an injection; we instead replay two of them without a model to obtain tool outputs with certain labels and evaluate the detector that screens those outputs.

\subsection{Detection-based defenses}
Detection-based defenses decide whether untrusted text contains an injection; the most widely deployed are fine-tuned encoders. ProtectAI's DeBERTa classifier (v2) is the detector baseline that AgentDojo ships and the model behind ProtectAI's LLM Guard scanner; PIGuard~\cite{piguard} (released earlier as InjecGuard) adds training data with trigger words to reduce over-defense; Meta's Prompt Guard models ship with LlamaFirewall~\cite{llamafirewall} as its scanner for untrusted input. A 2026 generation of open detectors, which we include, targets agent inputs explicitly: some are trained on documents and tool responses with planted instructions, and one adds role and source tokens to its input. These newer models are documented only by their model cards, so we audit what can be audited and report card claims as claims (\ref{app:configs}). Other detectors use the protected model itself. PromptShield~\cite{promptshield} trains a detector for deployment at low false-positive rates; DataSentinel~\cite{datasentinel} fine-tunes a detection LLM in a minimax game against adaptive attacks; Attention Tracker~\cite{attentiontracker} and Rennervate~\cite{rennervate} read attention patterns of the target model; and AlignSentinel~\cite{alignsentinel} separates instructions that align with the user's task from those that do not. Chen et al.~\cite{chen2025detected} evaluate LLMs and open-source detectors on indirect injections and remove detected instructions by segmentation or extraction. Our task-aware judges follow the intuition of AlignSentinel and Task Shield~\cite{taskshield} without training, and we use them only as a diagnostic, not as a proposed defense.

\subsection{Model-level, system-level, and task-aware defenses}
A second family does not try to recognize injections. StruQ~\cite{struq} and SecAlign~\cite{secalign} fine-tune the model to follow instructions only from the prompt channel, and the instruction hierarchy~\cite{instructionhierarchy} trains models to rank instructions by their source. Spotlighting~\cite{spotlighting} marks untrusted data with delimiters or encodings so the model can tell it apart. System-level defenses change the agent's architecture: CaMeL~\cite{camel} lets a privileged model plan from the user's request alone and tracks the provenance of every value derived from untrusted data, and MELON~\cite{melon} re-runs the agent with the user's request masked and flags an attack when the agent's tool calls do not depend on the request. Task Shield~\cite{taskshield} checks that every action serves the user's goal. These defenses do not depend on recognizing injected text, so our results about detectors do not transfer to them. Several of them report the false-positive cost of the detector baseline they compare against (Section~\ref{sec:community}), and a deployment that cannot adopt them, for instance because it uses a closed model or an existing agent framework, is exactly the deployment that falls back on a detector.

\subsection{Evaluation methodology, benchmark leakage, and pitfalls}
Our study belongs to a line of work that re-examines how security defenses are evaluated. Arp et al.~\cite{arp2022dos} catalog pitfalls such as data snooping and inappropriate baselines in machine learning for security, Pendlebury et al.~\cite{tesseract} show that malware classifiers look far better when training and test data are not separated in time, and Evertz et al.~\cite{chasingshadows} extend such a catalog to LLM security research. Sainz et al.~\cite{sainz2023} call for a contamination check for each LLM benchmark. For prompt injection, Nasr et al.~\cite{attackermovessecond} and Zhan et al.~\cite{zhan2025adaptive} show that defenses evaluated against static attacks fall to adaptive ones, and Jia et al.~\cite{jia2025critical} argue for evaluation against diverse and adaptive attacks and for effects on utility.

Several studies compare detectors directly. Dzhaliuk et al.~\cite{dzhaliuk2026} compare four classes of detection methods on more than 300{,}000 labeled prompts and find that fine-tuned encoders achieve the highest detection rate on balanced benchmarks. PIDS-Bench~\cite{pidsbench} evaluates detectors at fixed thresholds under over-defense, obfuscation, and distribution shift, and finds that a detector with F1 above 0.98 on its held-out split still misclassifies a third of an external benign set. Zizzo et al.~\cite{zizzo2024} find that simple baselines can match state-of-the-art guardrails out of distribution, Hackett et al.~\cite{hackett2025} evade detection systems with character injection and adversarial machine-learning attacks, and WAInjectBench~\cite{wainjectbench} finds that detectors for web agents miss attacks without explicit instructions. Fomin~\cite{benchmarkslie} shows with leave-one-dataset-out evaluation of activation probes that same-source splits inflate AUROC. Biswas~\cite{biswas2026} freezes a threshold calibrated on one benchmark, moves it across attack shifts, and finds that ProtectAI v2 and Prompt Guard~2 miss whole attack categories with near-certain confidence; that study deduplicates its evaluation sources against each other, but does not compare them with the detectors' training data. Khan et al.~\cite{khan2026} evaluate defenses, including a detector, on AgentDojo's banking suite with three commercial models and report security, utility, and run-to-run variation. Closest to our setting, the community benchmark buried-injections~\cite{buriedinjections} scores ten open detectors on AgentDojo injections embedded in tool outputs and recalibrates thresholds per suite; we compare its results with ours in Section~\ref{sec:community}.

These studies share our concern that headline accuracy misleads, and three of them (PIDS-Bench, Biswas, and buried-injections) also fix a low false-alarm budget. They differ from ours in what they measure. Most evaluate prompts or documents rather than the serialized outputs of an agent's tools; the two that use AgentDojo tool outputs study one benchmark; none audits a released detector's published training data against the benchmark used to rank it; and none asks whether a ranking on one benchmark predicts the ranking on another. The conclusion of Dzhaliuk et al.\ that fine-tuned models detect best on balanced benchmarks is the kind of result our study cautions against generalizing: the detector that ranks first on BIPIA, a fine-tuned model, ranks \BestBipiaRankOnAD of \NDetectors on AgentDojo at 1\% FPR.
